\documentclass[11pt,a4paper]{article}
\usepackage[utf8]{inputenc}
\usepackage[T1]{fontenc}
\usepackage[czech]{babel}
\usepackage{amsmath,amsthm,amssymb,amsfonts,mathtools}
\usepackage{geometry}
\usepackage{enumitem}
\usepackage{hyperref}
\geometry{margin=1in}

\title{Transport--Spectral Contradiction Program for the Riemann Hypothesis\\
\large Verze 26 -- full lock-form manuscript with operator-valued residue-side addendum}
\author{Archiv Všeho -- Imprint}
\date{\today}

\newtheorem{theorem}{Věta}[section]
\newtheorem{lemma}[theorem]{Lemma}
\newtheorem{corollary}[theorem]{Důsledek}
\newtheorem{definition}[theorem]{Definice}
\newtheorem{remark}[theorem]{Poznámka}
\newtheorem{proposition}[theorem]{Propozice}

\begin{document}

\maketitle

\begin{abstract}
Tento dokument podává plnou syntézu transportně-spektrálního programu pro Riemannovu hypotézu.
Strana \emph{uzavření} (closure) dává nulový excess přes Exact Window Law.
Strana \emph{obstrukce} (obstruction) říká, že každý mimokritický nulový bod by po transportu
vygeneroval lineární drift, jehož kvartická penalizace je koercivně neomezená.
Mostní nerovnost (bridge inequality) pak převádí tento drift do witness bank.
Výsledkem je čistý spor
\[
E_{\mathrm{ex}}(f)=0
\qquad\text{versus}\qquad
E_{\mathrm{ex}}(f)=+\infty.
\]
Dokument je matematicky poctivý: finální spor je zcela čistý, ale přechod od transportovaného
módu k úplné mostní nerovnosti je zatím formulován jako podmíněný bridge package.
\end{abstract}

\section{Základní definice}

\begin{definition}[Dipólově regularizovaný reziduální kanál]
Nechť \(\mu\) označuje explicitně-formulový reziduální distribuční kanál ve fixní
Guinand--Weilově normalizaci. Zvolme \(\psi\in\mathscr{S}(\mathbb{R})\) tak, že
\[
\widehat{\psi}(0)=0.
\]
Definujeme dipólově regularizovaný kanál
\[
f:=\psi * \mu.
\]
\end{definition}

\begin{definition}[Witness bank]
Zvolme dyadické škály
\[
a_j=2^j a_0,\qquad j=0,\dots,J,
\]
a kladné váhy \(w_j>0\). Pro funkci \(h\) definujeme druhou diferenci svědka
\[
(W_a h)(\tau):=h(\tau+a)-2h(\tau)+h(\tau-a).
\]
Je-li \(\varphi_{\tau,L}\) nezáporné okno lokalizované na \([\tau,\tau+L]\), pak definujeme bankovní energii
\[
Q_{\mathrm{bank}}(h;\tau,L):=
\sum_{j=0}^{J} w_j^2 \int_{\tau}^{\tau+L} |W_{a_j}h(u)|^2\,\varphi_{\tau,L}(u)\,du.
\]
V zjednodušeném zápisu, když \(\varphi_{\tau,L}\equiv 1\) na nosiči integrace, budeme psát
\[
Q_{\mathrm{bank}}(h;\tau,L)\asymp
\sum_{j=0}^{J} w_j^2 \int_{\tau}^{\tau+L} |W_{a_j}h(u)|^2\,du.
\]
\end{definition}

\begin{definition}[Excess functional]
Nechť \(E_0\ge 0\) je baseline density. Definujeme excess
\[
E_{\mathrm{ex}}(f):=
\limsup_{L\to\infty}\sup_{\tau\in\mathbb{R}}
\frac{Q_{\mathrm{bank}}(f;\tau,L)-E_0L}{L}.
\]
\end{definition}

\section{Strana A: closure}

\begin{definition}[Exact Window Law]
Řekneme, že kanál \(f\) splňuje přesný okenní zákon, pokud existují
\[
J_f:\mathbb{R}\to\mathbb{R}
\quad\text{omezená},\qquad
R(\tau,L),
\]
takové, že
\[
Q_{\mathrm{bank}}(f;\tau,L)
=
E_0L+J_f(\tau+L)-J_f(\tau)+R(\tau,L),
\]
a současně
\[
\sup_{\tau\in\mathbb{R}}|J_f(\tau)|<\infty,
\qquad
\lim_{L\to\infty}\sup_{\tau\in\mathbb{R}}\frac{|R(\tau,L)|}{L}=0.
\]
\end{definition}

\begin{theorem}[Vanishing of excess]
Jestliže \(f\) splňuje Exact Window Law, pak
\[
E_{\mathrm{ex}}(f)=0.
\]
\end{theorem}

\begin{proof}
Po vydělení délkou \(L\) dostaneme
\[
\frac{Q_{\mathrm{bank}}(f;\tau,L)-E_0L}{L}
=
\frac{J_f(\tau+L)-J_f(\tau)}{L}
+
\frac{R(\tau,L)}{L}.
\]
Vezmeme supremum přes \(\tau\):
\[
\sup_\tau\left|
\frac{Q_{\mathrm{bank}}(f;\tau,L)-E_0L}{L}
\right|
\le
\sup_\tau\left|
\frac{J_f(\tau+L)-J_f(\tau)}{L}
\right|
+
\sup_\tau \frac{|R(\tau,L)|}{L}.
\]
Z omezenosti \(J_f\) plyne
\[
\sup_\tau\left|
\frac{J_f(\tau+L)-J_f(\tau)}{L}
\right|
\le
\frac{2\sup_\tau |J_f(\tau)|}{L}\to 0,
\]
a z předpokladu i
\[
\sup_\tau \frac{|R(\tau,L)|}{L}\to 0.
\]
Tedy
\[
\limsup_{L\to\infty}\sup_\tau
\frac{Q_{\mathrm{bank}}(f;\tau,L)-E_0L}{L}=0,
\]
což je přesně
\[
E_{\mathrm{ex}}(f)=0.
\]
\end{proof}

\section{Strana B: transportní obstrukce}

\begin{definition}[Lifted drift]
Nechť
\[
\rho=\beta+i\gamma
\]
je netriviální nulový bod se \(\beta\neq \frac12\). V logaritmickém čase \(\tau=\log t\) definujeme
transportovaný drift
\[
\delta_\rho(\tau)=\frac{\beta-\tfrac12}{2\pi}\tau+r_\rho(\tau),
\]
kde \(r_\rho\) je omezený remainder.
\end{definition}

\begin{definition}[Kvartická penalizace]
Definujeme
\[
P_4(\rho;\tau,L):=\int_\tau^{\tau+L} |\delta_\rho(u)|^4\,du.
\]
\end{definition}

\begin{lemma}[Koercivní růst]
Jestliže \(\beta\neq \frac12\) a \(r_\rho\) je omezený, pak existují konstanty \(c,C>0\) tak, že pro
dost velká \(\tau\) a všechna \(L\ge 1\) platí
\[
P_4(\rho;\tau,L)\ge c\left|\beta-\frac12\right|^4\tau^4L-C\tau^3L.
\]
\end{lemma}

\begin{proof}
Položme
\[
a:=\frac{\beta-\tfrac12}{2\pi},\qquad \delta_\rho(u)=au+b(u),
\]
kde \(b\) je omezená. Rozvoj \(|au+b(u)|^4\) dává vedoucí člen \(a^4u^4\) a nižší členy nejvýše
řádu \(u^3\). Po integraci přes \([\tau,\tau+L]\) dostaneme tvrzení.
\end{proof}

\section{Bridge algebra}

\begin{definition}[Izolovaný transportovaný mód]
Předpokládejme splitting
\[
f=f_\rho+f_{\mathrm{rest}},
\]
kde \(f_\rho\) je mód asociovaný s mimokritickému nulovému bodu \(\rho\), zatímco
\(f_{\mathrm{rest}}\) obsahuje zbytek spektrálního obsahu.
\end{definition}

\begin{definition}[Witness coercivity]
Řekneme, že witness bank čte izolovaný mód koercivně, pokud existují konstanty
\[
\kappa>0,\qquad C_0\ge 0
\]
tak, že pro dost velká \(\tau\) platí
\[
\sum_{j=0}^{J} w_j^2 |W_{a_j}f_\rho(\tau)|^2
\ge
\kappa |\delta_\rho(\tau)|^4-C_0.
\]
\end{definition}

\begin{definition}[Sub-coercive leakage / no-hiding]
Řekneme, že remainder kanál je bank-subcritical, pokud existují
\[
\eta\in[0,\kappa),\qquad C_1\ge 0
\]
tak, že pro dost velká \(\tau\) platí
\[
\sum_{j=0}^{J} w_j^2
\Bigl(
2\Re\bigl(W_{a_j}f_\rho(\tau)\overline{W_{a_j}f_{\mathrm{rest}}(\tau)}\bigr)
+
|W_{a_j}f_{\mathrm{rest}}(\tau)|^2
\Bigr)
\ge
-\eta |\delta_\rho(\tau)|^4-C_1.
\]
\end{definition}

\begin{lemma}[Kvadratický rozklad banky]
Z rozkladu
\[
f=f_\rho+f_{\mathrm{rest}}
\]
plyne pro každý scale \(a_j\)
\[
|W_{a_j}f|^2
=
|W_{a_j}f_\rho|^2
+
2\Re\!\Bigl(W_{a_j}f_\rho\,\overline{W_{a_j}f_{\mathrm{rest}}}\Bigr)
+
|W_{a_j}f_{\mathrm{rest}}|^2.
\]
Po sečtení přes \(j\) a integraci přes okno tedy
\[
Q_{\mathrm{bank}}(f;\tau,L)
=
Q_{\mathrm{bank}}(f_\rho;\tau,L)
+
Q_{\mathrm{cross}}(\rho,\mathrm{rest};\tau,L)
+
Q_{\mathrm{bank}}(f_{\mathrm{rest}};\tau,L).
\]
\end{lemma}

\begin{proof}
Jde o čistý kvadratický rozvoj výrazu \(|W_{a_j}(f_\rho+f_{\mathrm{rest}})|^2\).
\end{proof}

\begin{theorem}[Target Bridge Inequality]
Za předpokladů witness coercivity a no-hiding leakage existují konstanty
\[
c_1>0,\qquad c_2\ge 0
\]
tak, že pro dost velká \(\tau\) a všechna \(L\ge 1\) platí
\[
Q_{\mathrm{bank}}(f;\tau,L)-E_0L
\ge
c_1P_4(\rho;\tau,L)-c_2L.
\]
Přesněji lze vzít
\[
c_1:=\kappa-\eta>0,
\qquad
c_2:=C_0+C_1.
\]
\end{theorem}

\begin{proof}
Z koercivity dostáváme po integraci přes \([\tau,\tau+L]\)
\[
Q_{\mathrm{bank}}(f_\rho;\tau,L)\ge \kappa P_4(\rho;\tau,L)-C_0L.
\]
Z no-hiding estimate dostáváme
\[
Q_{\mathrm{cross}}(\rho,\mathrm{rest};\tau,L)
+
Q_{\mathrm{bank}}(f_{\mathrm{rest}};\tau,L)
\ge
-\eta P_4(\rho;\tau,L)-C_1L.
\]
Po sečtení:
\[
Q_{\mathrm{bank}}(f;\tau,L)-E_0L
\ge
(\kappa-\eta)P_4(\rho;\tau,L)-(C_0+C_1)L.
\]
To je přesně požadovaná nerovnost.
\end{proof}


\begin{lemma}[M5: Asymptotická ortogonalita driftu a no-hiding estimate]
Nechť
\[
f=f_\rho+f_{\mathrm{rest}}
\]
je splitting kanálu na izolovaný mimokritický mód a remainder.
Předpokládejme, že pro každé měřítko \(a_j\), \(j=0,\dots,J\), existují na okně
\([\tau,\tau+L]\) reprezentace
\[
W_{a_j}f_\rho(u)=B_{\rho,j}(u)e^{i\Theta_\rho(u)},
\qquad
W_{a_j}f_{\mathrm{rest}}(u)=B_{\mathrm{rest},j}(u)e^{i\Theta_{\mathrm{rest},j}(u)},
\]
kde platí následující podmínky:

\begin{enumerate}[label=\textbf{(\arabic*)}]
    \item \textbf{Fázová separace.}
    Pro rozdílovou fázi
    \[
    \Phi_j(u):=\Theta_\rho(u)-\Theta_{\mathrm{rest},j}(u)
    \]
    existují konstanty \(c_\Phi>0\), \(C_\Phi\ge 0\) tak, že pro všechna dost velká \(\tau\),
    všechna \(u\in[\tau,\tau+L]\) a všechna \(j\) platí
    \[
    |\Phi_j'(u)|\ge c_\Phi \tau,
    \qquad
    |\Phi_j''(u)|\le C_\Phi.
    \]

    \item \textbf{Kontrola amplitud.}
    Existují konstanty \(C_\rho,C_{\mathrm{rest}}\ge 0\) tak, že pro všechna dost velká \(\tau\),
    všechna \(u\in[\tau,\tau+L]\) a všechna \(j\) platí
    \[
    |B_{\rho,j}(u)|+|B_{\rho,j}'(u)|\le C_\rho \tau^2,
    \]
    \[
    |B_{\mathrm{rest},j}(u)|+|B_{\mathrm{rest},j}'(u)|\le C_{\mathrm{rest}}.
    \]

    \item \textbf{Koercivní dolní mez pro drift.}
    Existuje konstanta \(c_4>0\) tak, že pro všechna dost velká \(\tau\) a všechna \(L\ge 1\) platí
    \[
    P_4(\rho;\tau,L)\ge c_4\tau^4L.
    \]
\end{enumerate}

Pak existuje konstanta \(C_*>0\), nezávislá na \(\tau\) a \(L\), taková, že
\[
|Q_{\mathrm{cross}}(\rho,\mathrm{rest};\tau,L)|\le C_* \tau L.
\]
Dále tedy platí
\[
Q_{\mathrm{cross}}(\rho,\mathrm{rest};\tau,L)
+
Q_{\mathrm{bank}}(f_{\mathrm{rest}};\tau,L)
\ge
-\eta(\tau)\,P_4(\rho;\tau,L),
\]
kde
\[
\eta(\tau):=\frac{C_*}{c_4\tau^3}\xrightarrow[\tau\to\infty]{}0.
\]
Zvláště: je-li \(\kappa>0\) konstanta witness coercivity z M4, pak existuje \(\tau_*\) tak, že pro
všechna \(\tau\ge \tau_*\) platí
\[
\eta(\tau)<\frac{\kappa}{2},
\]
a tedy
\[
Q_{\mathrm{cross}}(\rho,\mathrm{rest};\tau,L)
+
Q_{\mathrm{bank}}(f_{\mathrm{rest}};\tau,L)
\ge
-\frac{\kappa}{2}\,P_4(\rho;\tau,L).
\]
\end{lemma}

\begin{proof}
Pro každé \(j\) definujme křížový integrál
\[
I_j(\tau,L):=
\int_\tau^{\tau+L}
W_{a_j}f_\rho(u)\,\overline{W_{a_j}f_{\mathrm{rest}}(u)}\,du.
\]
Po dosazení fázově-amplitudové reprezentace dostáváme
\[
I_j(\tau,L)=\int_\tau^{\tau+L} g_j(u)e^{i\Phi_j(u)}\,du,
\]
kde
\[
g_j(u):=B_{\rho,j}(u)\,\overline{B_{\mathrm{rest},j}(u)}.
\]

Z předpokladů o amplitudách plyne, že existuje konstanta \(C_g>0\) tak, že
\[
|g_j(u)|\le C_g \tau^2,
\qquad
|g_j'(u)|\le C_g \tau^2
\]
na celém intervalu \([\tau,\tau+L]\), pro všechna dost velká \(\tau\) a všechna \(j\).

Použijeme integraci po partiích v nestacionární fázi:
\[
I_j(\tau,L)
=
\int_\tau^{\tau+L}
\frac{g_j(u)}{i\Phi_j'(u)}
\frac{d}{du}\bigl(e^{i\Phi_j(u)}\bigr)\,du.
\]
Tedy
\[
I_j(\tau,L)
=
\left[
\frac{g_j(u)}{i\Phi_j'(u)}e^{i\Phi_j(u)}
\right]_{\tau}^{\tau+L}
-
\int_\tau^{\tau+L}
\frac{d}{du}\left(\frac{g_j(u)}{i\Phi_j'(u)}\right)e^{i\Phi_j(u)}\,du.
\]
Odtud dostaneme odhad
\[
|I_j(\tau,L)|
\le
\frac{2\sup_{u\in[\tau,\tau+L]}|g_j(u)|}{c_\Phi\tau}
+
\int_\tau^{\tau+L}
\left(
\frac{|g_j'(u)|}{c_\Phi\tau}
+
\frac{|g_j(u)|\,|\Phi_j''(u)|}{c_\Phi^2\tau^2}
\right)du.
\]
Po dosazení odhadů
\[
|g_j(u)|\le C_g\tau^2,\qquad |g_j'(u)|\le C_g\tau^2,\qquad |\Phi_j''(u)|\le C_\Phi
\]
dostáváme
\[
|I_j(\tau,L)|
\le
C_1\tau + C_2L\tau + C_3L.
\]
Protože \(L\ge 1\), můžeme absorbovat okrajový člen do lineárního odhadu a dostáváme
\[
|I_j(\tau,L)|\le C_j^*\,L\tau
\]
pro vhodnou konstantu \(C_j^*>0\).

Křížový člen banky tedy splňuje
\[
|Q_{\mathrm{cross}}(\rho,\mathrm{rest};\tau,L)|
\le
2\sum_{j=0}^J w_j^2 |I_j(\tau,L)|
\le
C_*\,L\tau
\]
pro vhodnou konstantu \(C_*>0\).

Dále, protože
\[
Q_{\mathrm{bank}}(f_{\mathrm{rest}};\tau,L)\ge 0,
\]
máme
\[
Q_{\mathrm{cross}}(\rho,\mathrm{rest};\tau,L)
+
Q_{\mathrm{bank}}(f_{\mathrm{rest}};\tau,L)
\ge
-|Q_{\mathrm{cross}}(\rho,\mathrm{rest};\tau,L)|
\ge
-C_*L\tau.
\]
Použijeme koercivní dolní mez
\[
P_4(\rho;\tau,L)\ge c_4\tau^4L.
\]
Pak
\[
C_*L\tau
=
\frac{C_*}{c_4\tau^3}\,c_4\tau^4L
\le
\frac{C_*}{c_4\tau^3}\,P_4(\rho;\tau,L).
\]
Definujeme tedy
\[
\eta(\tau):=\frac{C_*}{c_4\tau^3}.
\]
Tím dostáváme
\[
Q_{\mathrm{cross}}(\rho,\mathrm{rest};\tau,L)
+
Q_{\mathrm{bank}}(f_{\mathrm{rest}};\tau,L)
\ge
-\eta(\tau)\,P_4(\rho;\tau,L),
\]
a protože \(\eta(\tau)\to 0\), existuje \(\tau_*\) tak, že pro \(\tau\ge \tau_*\) platí
\[
\eta(\tau)<\frac{\kappa}{2}.
\]
Tedy
\[
Q_{\mathrm{cross}}(\rho,\mathrm{rest};\tau,L)
+
Q_{\mathrm{bank}}(f_{\mathrm{rest}};\tau,L)
\ge
-\frac{\kappa}{2}\,P_4(\rho;\tau,L).
\]
Důkaz je hotov.
\end{proof}

\begin{corollary}[M5 implies no-hiding]
Za předpokladů witness coercivity a předchozího lemmatu existují konstanty
\[
c_1>0,\qquad c_2\ge 0
\]
tak, že pro všechna dost velká \(\tau\) a všechna \(L\ge 1\) platí
\[
Q_{\mathrm{bank}}(f;\tau,L)-E_0L
\ge
c_1P_4(\rho;\tau,L)-c_2L.
\]
Konkrétně lze vzít
\[
c_1=\frac{\kappa}{2}.
\]
\end{corollary}

\begin{proof}
Witness coercivity dává po integraci přes \([\tau,\tau+L]\)
\[
Q_{\mathrm{bank}}(f_\rho;\tau,L)\ge \kappa P_4(\rho;\tau,L)-C_0L.
\]
Předchozí lemma dává
\[
Q_{\mathrm{cross}}(\rho,\mathrm{rest};\tau,L)
+
Q_{\mathrm{bank}}(f_{\mathrm{rest}};\tau,L)
\ge
-\frac{\kappa}{2}P_4(\rho;\tau,L).
\]
Po sečtení obou odhadů dostaneme
\[
Q_{\mathrm{bank}}(f;\tau,L)-E_0L
\ge
\frac{\kappa}{2}P_4(\rho;\tau,L)-C_0L.
\]
To je požadovaná target bridge inequality.
\end{proof}



\section{Explosive obstruction}

\begin{theorem}[Pozitivní excess z mimokritického bodu]
Jestliže existuje off-critical nulový bod \(\rho=\beta+i\gamma\), \(\beta\neq \frac12\), a platí
Target Bridge Inequality, pak
\[
E_{\mathrm{ex}}(f)>0.
\]
\end{theorem}

\begin{proof}
Po vydělení délkou \(L\) máme
\[
\frac{Q_{\mathrm{bank}}(f;\tau,L)-E_0L}{L}
\ge
c_1\frac{P_4(\rho;\tau,L)}{L}-c_2.
\]
Použitím koercivního odhadu dostáváme
\[
\frac{Q_{\mathrm{bank}}(f;\tau,L)-E_0L}{L}
\ge
c_1\Bigl(c\left|\beta-\frac12\right|^4\tau^4-C\tau^3\Bigr)-c_2.
\]
Pro dost velká \(\tau\) je pravá strana kladná. Vezmeme \(\sup_\tau\) a poté \(\limsup_{L\to\infty}\),
odkud
\[
E_{\mathrm{ex}}(f)>0.
\]
\end{proof}

\begin{corollary}[Bridge Lock]
Jestliže je Target Bridge Inequality splněna s \(c_1>0\), pak kvartický drift nelze dlouhodobě
skrýt: pro dost velká \(\tau\) a všechna \(L\ge 1\) je
\[
\frac{Q_{\mathrm{bank}}(f;\tau,L)-E_0L}{L}>0.
\]
\end{corollary}

\section{Finální spor}

\begin{theorem}[Conditional Contradiction Theorem]
Nechť \(f=\psi*\mu\). Předpokládejme:
\begin{enumerate}[label=\textbf{(\arabic*)}]
\item \(f\) splňuje Exact Window Law, tedy \(E_{\mathrm{ex}}(f)=0\);
\item každý mimokritický nulový bod \(\rho=\beta+i\gamma\), \(\beta\neq \frac12\), generuje
transportovaný mód \(f_\rho\), pro nějž platí witness coercivity a no-hiding leakage, a tedy
Target Bridge Inequality.
\end{enumerate}
Pak platí Riemannova hypotéza.
\end{theorem}

\begin{proof}
Z bodu \textbf{(1)} plyne
\[
E_{\mathrm{ex}}(f)=0.
\]
Předpokládejme pro spor existenci off-critical nulového bodu
\[
\rho=\beta+i\gamma,\qquad \beta\neq \frac12.
\]
Pak z bodu \textbf{(2)} a předchozí věty plyne
\[
E_{\mathrm{ex}}(f)>0.
\]
To je spor s \(E_{\mathrm{ex}}(f)=0\). Proto žádný mimokritický nulový bod neexistuje. Každý
netriviální nulový bod tedy musí ležet na kritické přímce:
\[
\Re(s)=\frac12.
\]
\end{proof}

\section{Co přesně ještě zbývá}

Aby se předchozí věta stala bezpodmínečnou, zbývá analyticky uzavřít přesně bridge package.
V pracovním jazyce to znamená:

\begin{enumerate}[label=\textbf{M\arabic*.}]
\item explicitní single-zero extrakce \(f_\rho\) ve správné transportní normalizaci;
\item parity survival netriviální off-critical složky;
\item odvození transportního zákona
\[
\delta_\rho(\tau)=\frac{\beta-\tfrac12}{2\pi}\tau+r_\rho(\tau),
\qquad r_\rho \text{ bounded};
\]
\item skutečný důkaz witness coercivity;
\item bezpodmínečné ověření předpokladů M5, zejména fázové separace, amplitudové kontroly a z nich plynoucí nerovnosti
\[
\eta(\tau)<\kappa \text{ pro dost velká }\tau.
\]
\end{enumerate}

\begin{remark}[Nejcitlivější zadrhel]
Finální algebraická syntéza je čistá. Kritický analytický bod je potvrdit
\[
\eta<\kappa
\]
v no-hiding estimate tak, aby remainder channel nikdy nepřebil hlavní koercivní driftový příspěvek.
Právě tady se podmíněné tvrzení mění v absolutní fakt.
\end{remark}


\section{Catalan parity route: strukturální upgrade pro M4/M5}

V této sekci přidáváme Catalanovu vrstvu jako \emph{strukturální šablonu} pro možný
bezpodmínečný upgrade modulů M4 a M5. Tato vrstva nevstupuje do dokumentu jako již uzavřený
důkaz RH, ale jako přesně formulovaná cesta, která by odstranila závislost na externích
parametrech \(\kappa\) a \(\eta\), pokud se podaří analyticky propojit paritní oscilátor s reálným
transportovaným módem \(f_\rho\).

\subsection*{A. Co je v Catalanově modulu skutečně dokázáno}

\begin{definition}[Paritní oscilátor]
Definujme diskrétní dynamický systém
\[
x_{n+1}=-x_n,\qquad x_0=1.
\]
\end{definition}

\begin{lemma}[Řešení paritního oscilátoru]
Platí
\[
x_n=(-1)^n.
\]
\end{lemma}

\begin{proof}
Důkaz je triviální indukcí.
\end{proof}

\begin{definition}[Absolutní energie]
Definujme
\[
E_n:=|x_n|.
\]
\end{definition}

\begin{theorem}[Invariance energie]
Pro všechna \(n\in\mathbb{N}\) platí
\[
E_n=1.
\]
\end{theorem}

\begin{proof}
Protože \(|(-1)^n|=1\) pro všechna \(n\), tvrzení je okamžité.
\end{proof}

\begin{definition}[Odd-parity projection]
Pro posloupnost \((a_n)\) definujme
\[
(P_{\mathrm{odd}}a)_n:=\frac12\bigl(1-(-1)^n\bigr)a_n.
\]
\end{definition}

\begin{lemma}[Vlastnosti operátoru \(P_{\mathrm{odd}}\)]
Platí:
\begin{enumerate}[label=\textbf{(\arabic*)}]
\item \(P_{\mathrm{odd}}^2=P_{\mathrm{odd}}\),
\item \((P_{\mathrm{odd}}a)_n=0\) pro sudá \(n\),
\item \((P_{\mathrm{odd}}a)_n=a_n\) pro lichá \(n\).
\end{enumerate}
\end{lemma}

\begin{proof}
Stačí použít identitu
\[
\bigl(1-(-1)^n\bigr)^2=2\bigl(1-(-1)^n\bigr).
\]
\end{proof}

\begin{theorem}[Catalan--mini-zeta strukturální identita]
Nechť
\[
I=\int_0^1\int_0^1 \frac{\log(1+xy)}{1+x^2}\,dx\,dy,
\qquad
J=\int_0^{\pi/4}\log(1+\tan\theta)\,d\theta.
\]
Nechť \(G\) je Catalanova konstanta a
\[
S_{\mathrm{raw}}=
\sum_{n=1}^\infty \frac{(-1)^{n+1}}{n^2}\log\!\left(1+\frac1n\right).
\]
Pak platí přesná strukturální identita
\[
I=2J-\frac14 G-\frac14 S_{\mathrm{raw}}.
\]
\end{theorem}

\begin{remark}
Catalanův blok tedy opravdu dodává tři pevné věci:
\begin{enumerate}[label=\textbf{(\arabic*)}]
\item invariantní absolutní energii paritního oscilátoru,
\item exaktní odd-parity filtr \(P_{\mathrm{odd}}\),
\item strukturální vazbu mezi paritním kanálem, Catalanovou konstantou a mini-zeta Dirichletovým kanálem.
\end{enumerate}
To jsou skutečná tvrzení.
\end{remark}

\subsection*{B. Navržený bezpodmínečný upgrade M4}

Motivace je tato: v předchozí části dokumentu je nejcitlivější místo M4 otázka, zda amplituda
transportovaného módu
\[
f_\rho(\tau)=A_\rho(\tau)e^{i\Theta_\rho(\tau)}
\]
nemůže asymptoticky kolabovat k nule. Catalanův modul nabízí následující strukturální cestu:

\begin{proposition}[Navržená parity-coercivity route]
Předpokládejme, že izolovaný mód \(f_\rho\) lze přirozeně realizovat v paritním sektoru,
tj. že existuje reprezentace, ve které
\[
f_\rho \longmapsto P_{\mathrm{odd}}(f_\rho),
\]
a současně energie tohoto sektoru dědí invariant absolutní hodnoty paritního oscilátoru.
Pak existuje strukturální dolní mez
\[
A_\rho(\tau)\ge A_0>0,
\]
a tedy witness coercivity z M4 lze psát ve tvaru
\[
\sum_{j=0}^{J} w_j^2 |W_{a_j}f_\rho(\tau)|^2
\ge
\kappa |\delta_\rho(\tau)|^4-C_0
\]
s \(\kappa>0\) nezávislou na \(\tau\).
\end{proposition}

\begin{remark}[Poctivý status]
Tohle zatím není hotová věta vyplývající pouze z Catalanova modulu. Je to přesně formulovaná
\emph{route}: aby se z ní stal důkaz, musí být analyticky ukázáno, že reálný transportovaný mód
RH skutečně leží v takovém paritním sektoru a že jeho amplitudová energie je invariantem
přenositelným z diskrétního oscilátoru na witness-level čtení.
\end{remark}

\subsection*{C. Navržený bezpodmínečný upgrade M5}

V M5 je nejcitlivější bod nerovnost
\[
\eta(\tau)<\kappa
\]
pro dost velká \(\tau\). Catalanova vrstva navrhuje silnější variantu: místo pouhé asymptotické
malosti úniku se pokusit remainder kanál \emph{strukturálně vynulovat} v paritním sektoru.

\begin{proposition}[Navržený spectral firewall]
Předpokládejme, že:
\begin{enumerate}[label=\textbf{(\arabic*)}]
\item driftový mód \(f_\rho\) je nesen čistě Catalanovým paritním blokem;
\item remainder \(f_{\mathrm{rest}}\) nemá v tomto paritním sektoru vedoucí stopu;
\item strukturální identita
\[
I=2J-\frac14 G-\frac14 S_{\mathrm{raw}}
\]
je stabilní vůči projekci kanálu do witness bank.
\end{enumerate}
Pak lze očekávat, že v paritním sektoru platí exaktní eliminace interference:
\[
\eta\equiv 0.
\]
\end{proposition}

\begin{remark}[Poctivý status]
Ani toto zatím není hotový důkaz. Catalanova identita je strukturální, ale skok od ní k tvrzení
\[
\eta\equiv 0
\]
vyžaduje nový analytický článek: důkaz, že remainder kanál skutečně nemá vedoucí komponentu
v tomtéž paritním witness sektoru jako drift.
\end{remark}

\subsection*{D. Co by z toho plynulo, kdyby se route uzavřela}

Pokud by se podařilo rigorózně odvodit obě předchozí propozice, dostali bychom:
\[
\kappa>0
\qquad\text{a dokonce případně}\qquad
\eta=0.
\]
Pak by mostní nerovnost přestala být podmínkou a stala by se strukturálním důsledkem paritního
kanálu:
\[
Q_{\mathrm{bank}}(f;\tau,L)-E_0L
\ge
\kappa P_4(\rho;\tau,L)-C L.
\]
V kombinaci s closure stranou by to uzavřelo RH bezpodmínečně.

\begin{remark}[Archiv Všeho reading]
Catalanova vrstva tedy může být čtena jako kandidát na \emph{energetický invariant} a
\emph{paritní firewall} současně:
\begin{enumerate}[label=\textbf{(\arabic*)}]
\item invariance energie brání odpaření driftového módu,
\item paritní filtr eliminuje nebo asymptoticky vytlačuje remainder interference,
\item strukturální identita drží pohromadě Catalanův blok a mini-zeta kanál.
\end{enumerate}
To je přesně důvod, proč je tato vrstva cenná pro unconditional route.
\end{remark}

\subsection*{E. Finální poctivý verdict k Catalan route}

Catalanova vrstva výrazně posiluje dokument na úrovni \emph{struktury}:
\begin{enumerate}[label=\textbf{(\arabic*)}]
\item dává pevný paritní oscilátor s invariantní energií,
\item dává exaktní projekční filtr \(P_{\mathrm{odd}}\),
\item dává strukturální spojení s Catalanovou konstantou a mini-zeta kanálem.
\end{enumerate}
Ale sama o sobě ještě nedokazuje, že pro RH mód automaticky platí
\[
\kappa\ge G>0
\qquad\text{nebo}\qquad
\eta\equiv 0.
\]
Tyto výroky zůstávají \emph{návrhem bezpodmínečné cesty}, nikoli hotovým theoremem.



\section{Version 18 corrected: asymptotic firewall via nonstationary phase}

V této sekci formulujeme korektní analytickou verzi asymptotického firewallu.
Klíčová oprava je tato:
nebudeme předpokládat, že remainder kanál má globálně omezené frekvenční spektrum.
Místo toho budeme pracovat s \emph{rozdílovou fází} mezi driftovým módem a remainderem
na každém relevantním bank-scale a využijeme nestacionární fázi.

\subsection*{Fixace proměnné}

Od této chvíle pracujeme výhradně v transportní proměnné \(\tau\).
Všechny fáze, frekvence, bank windows i odhady jsou formulovány vzhledem k \(\tau\),
nikoli k původní proměnné \(t\).

\subsection*{Driftová fáze a remainder fáze}

Nechť
\[
f=f_\rho+f_{\mathrm{rest}}
\]
je splitting kanálu na izolovaný mimokritický mód a remainder.
Pro každý scale \(a_j\), \(j=0,\dots,J\), předpokládejme reprezentace
\[
W_{a_j}f_\rho(u)=B_{\rho,j}(u)e^{i\Theta_\rho(u)},
\qquad
W_{a_j}f_{\mathrm{rest}}(u)=B_{\mathrm{rest},j}(u)e^{i\Theta_{\mathrm{rest},j}(u)}
\]
na okně \([\tau,\tau+L]\).

Definujme rozdílovou fázi
\[
\Phi_j(u):=\Theta_\rho(u)-\Theta_{\mathrm{rest},j}(u).
\]

\subsection*{Frekvenční mezera v korektním tvaru}

\begin{definition}[Lokální frekvenční mezera]
Řekneme, že driftový mód je na okně \([\tau,\tau+L]\) frekvenčně separován od remainderu,
jestliže existuje konstanta \(c_\Phi>0\) tak, že pro všechna dost velká \(\tau\),
všechna \(u\in[\tau,\tau+L]\) a všechna \(j\) platí
\[
|\Phi_j'(u)|\ge c_\Phi \tau.
\]
\end{definition}

\begin{remark}
To je správná náhrada za dřívější příliš silné tvrzení, že remainder má pevně omezený
Fourierův support \([0,\Omega]\). Pro M5 ve skutečnosti nepotřebujeme bounded-band hypothesis;
stačí, aby rozdílová fáze měla na relevantním okně derivaci uniformně vzdálenou od nuly.
\end{remark}

\subsection*{Nestacionární-fázové lemma pro jeden scale}

\begin{lemma}[One-scale asymptotic orthogonality]
Předpokládejme, že pro pevné \(j\) platí na \([\tau,\tau+L]\):
\begin{enumerate}[label=\textbf{(\arabic*)}]
    \item \(\Phi_j\in C^2\) a
    \[
    |\Phi_j'(u)|\ge c_\Phi\tau,
    \qquad
    |\Phi_j''(u)|\le C_\Phi;
    \]
    \item amplitudy \(B_{\rho,j},B_{\mathrm{rest},j}\in C^1\) a existuje konstanta \(C_B>0\) tak, že
    \[
    |B_{\rho,j}(u)|+|B'_{\rho,j}(u)|\le C_B\tau^2,
    \]
    \[
    |B_{\mathrm{rest},j}(u)|+|B'_{\mathrm{rest},j}(u)|\le C_B.
    \]
\end{enumerate}
Pak existuje konstanta \(C_j^*>0\) tak, že
\[
\left|
\int_\tau^{\tau+L}
W_{a_j}f_\rho(u)\,\overline{W_{a_j}f_{\mathrm{rest}}(u)}\,du
\right|
\le
C_j^*\,L\tau.
\]
\end{lemma}

\begin{proof}
Položme
\[
g_j(u):=B_{\rho,j}(u)\overline{B_{\mathrm{rest},j}(u)}.
\]
Pak
\[
\int_\tau^{\tau+L}
W_{a_j}f_\rho(u)\,\overline{W_{a_j}f_{\mathrm{rest}}(u)}\,du
=
\int_\tau^{\tau+L} g_j(u)e^{i\Phi_j(u)}\,du.
\]
Z předpokladů plyne
\[
|g_j(u)|\le C_g\tau^2,
\qquad
|g_j'(u)|\le C_g\tau^2
\]
pro vhodnou konstantu \(C_g>0\).

Použijeme integraci po partiích v nestacionární fázi:
\[
\int_\tau^{\tau+L} g_j(u)e^{i\Phi_j(u)}\,du
=
\int_\tau^{\tau+L}
\frac{g_j(u)}{i\Phi_j'(u)}
\frac{d}{du}\bigl(e^{i\Phi_j(u)}\bigr)\,du.
\]
Tedy
\[
\int_\tau^{\tau+L} g_j(u)e^{i\Phi_j(u)}\,du
=
\left[
\frac{g_j(u)}{i\Phi_j'(u)}e^{i\Phi_j(u)}
\right]_{\tau}^{\tau+L}
-
\int_\tau^{\tau+L}
\frac{d}{du}\left(\frac{g_j(u)}{i\Phi_j'(u)}\right)e^{i\Phi_j(u)}\,du.
\]
Odtud
\[
\left|
\int_\tau^{\tau+L} g_j(u)e^{i\Phi_j(u)}\,du
\right|
\le
\frac{2\sup |g_j|}{c_\Phi\tau}
+
\int_\tau^{\tau+L}
\left(
\frac{|g_j'|}{c_\Phi\tau}
+
\frac{|g_j|\,|\Phi_j''|}{c_\Phi^2\tau^2}
\right)du.
\]
Po dosazení odhadů dostaneme
\[
\left|
\int_\tau^{\tau+L} g_j(u)e^{i\Phi_j(u)}\,du
\right|
\le
C_1\tau + C_2L\tau + C_3L.
\]
Protože \(L\ge 1\), můžeme psát
\[
\left|
\int_\tau^{\tau+L}
W_{a_j}f_\rho(u)\,\overline{W_{a_j}f_{\mathrm{rest}}(u)}\,du
\right|
\le
C_j^*\,L\tau.
\]
\end{proof}

\subsection*{Finite-bank asymptotic firewall}

\begin{lemma}[Finite-bank asymptotic firewall]
Za předpokladů předchozího lemmatu pro všechna \(j=0,\dots,J\) existuje konstanta \(C_*>0\) tak, že
\[
|Q_{\mathrm{cross}}(\rho,\mathrm{rest};\tau,L)|\le C_*L\tau.
\]
Je-li navíc pro dost velká \(\tau\) splněno
\[
P_4(\rho;\tau,L)\ge c_4\tau^4L
\]
s nějakou konstantou \(c_4>0\), pak
\[
\frac{|Q_{\mathrm{cross}}(\rho,\mathrm{rest};\tau,L)|}{P_4(\rho;\tau,L)}
\le
\frac{C_*}{c_4\tau^3}\xrightarrow[\tau\to\infty]{}0.
\]
\end{lemma}

\begin{proof}
Z definice křížového členu a z jednoskalového odhadu máme
\[
|Q_{\mathrm{cross}}(\rho,\mathrm{rest};\tau,L)|
\le
2\sum_{j=0}^{J} w_j^2
\left|
\int_\tau^{\tau+L}
W_{a_j}f_\rho(u)\,\overline{W_{a_j}f_{\mathrm{rest}}(u)}\,du
\right|.
\]
Každý člen je \(\le C_j^*L\tau\), tedy
\[
|Q_{\mathrm{cross}}(\rho,\mathrm{rest};\tau,L)|
\le
C_*L\tau
\]
pro vhodnou konstantu \(C_*>0\).
Po vydělení koercivním dolním odhadem pro \(P_4\) dostáváme
\[
\frac{|Q_{\mathrm{cross}}(\rho,\mathrm{rest};\tau,L)|}{P_4(\rho;\tau,L)}
\le
\frac{C_*}{c_4\tau^3}\to 0.
\]
\end{proof}

\begin{corollary}[Corrected M5 route]
Jestliže současně platí witness coercivity z M4
\[
Q_{\mathrm{bank}}(f_\rho;\tau,L)\ge \kappa P_4(\rho;\tau,L)-C_0L
\]
a finite-bank asymptotic firewall, pak existuje \(\tau_*\) tak, že pro všechna \(\tau\ge \tau_*\)
je remainder leakage subkoercivní:
\[
Q_{\mathrm{cross}}(\rho,\mathrm{rest};\tau,L)
=
o\!\left(P_4(\rho;\tau,L)\right).
\]
Zvláště existuje \(\eta(\tau)\to 0\) a od jistého okamžiku platí
\[
\eta(\tau)<\kappa.
\]
\end{corollary}

\begin{remark}
Tohle je správná theorem-ready forma v18. Stále ovšem zbývá bezpodmínečně doložit
předpoklady:
\begin{enumerate}[label=\textbf{(\arabic*)}]
    \item fázovou separaci \(|\Phi_j'(u)|\ge c_\Phi\tau\),
    \item amplitudovou kontrolu,
    \item kompatibilitu s reálnou residue-side normalizací RH witnesse.
\end{enumerate}
Takže v18 corrected je silná route pro M5, ale ještě ne finální unconditional uzávěr.
\end{remark}



\section{CROWN stability seal: phase rigidity as a structural stabilizer}

Tato sekce integruje modul \(\psi\mathrm{CROWN}\) jako \emph{stabilizační vrstvu} nad
transportním mostem. Její role je přesně vymezena:
\begin{enumerate}
    \item neposkytuje sama o sobě důkaz RH,
    \item ale dává čistý modelový theorem o tom, kdy diagonální evoluce zachovává
    plnou fázovou strukturu,
    \item a tím zpřesňuje, co přesně musí být v RH kanálu dokázáno, aby se z
    ``koherentního driftového módu'' stal skutečně rigidní witness stav.
\end{enumerate}

\subsection*{Definice crown stavu}

\begin{definition}[\(\psi\mathrm{CROWN}\) stav]
Nechť \(H\) je konečně-dimenzionální komplexní Hilbertův prostor s ortonormální bází
\[
\{|0\rangle,|1\rangle,\dots,|n\rangle\}.
\]
Pro fáze \(\theta_0,\dots,\theta_n\in\mathbb{R}\) definujeme stav
\[
\psi\mathrm{CROWN}
=
\frac{1}{\sqrt{n+1}}
\sum_{k=0}^{n} e^{i\theta_k}|k\rangle.
\]
\end{definition}

\begin{lemma}[Normalizace]
Platí
\[
\langle \psi\mathrm{CROWN},\psi\mathrm{CROWN}\rangle=1.
\]
\end{lemma}

\begin{proof}
Z ortonormality báze plyne
\[
\langle \psi\mathrm{CROWN},\psi\mathrm{CROWN}\rangle
=
\frac{1}{n+1}\sum_{j=0}^{n}\sum_{k=0}^{n}
e^{i(\theta_k-\theta_j)}\langle j|k\rangle
=
\frac{1}{n+1}\sum_{k=0}^{n}1=1.
\]
\end{proof}

\subsection*{Diagonální unitární evoluce}

\begin{definition}[Diagonální evoluce]
Nechť operátor \(H\) působí diagonálně:
\[
H|k\rangle=\lambda_k|k\rangle,
\qquad \lambda_k\in\mathbb{R}.
\]
Pak unitární evoluce
\[
U(t):=e^{-itH}
\]
splňuje
\[
U(t)|k\rangle=e^{-it\lambda_k}|k\rangle.
\]
\end{definition}

\begin{lemma}[Explicitní evoluce crown stavu]
Pro každé \(t\in\mathbb{R}\) platí
\[
U(t)\psi\mathrm{CROWN}
=
\frac{1}{\sqrt{n+1}}
\sum_{k=0}^{n} e^{i(\theta_k-\lambda_k t)}|k\rangle.
\]
\end{lemma}

\begin{proof}
Tvrzení plyne přímo z linearity a diagonálního působení \(U(t)\) na bázi.
\end{proof}

\begin{theorem}[Probability stability]
Pro každé \(k\in\{0,\dots,n\}\) a každé \(t\in\mathbb{R}\) platí
\[
|\langle k|U(t)\psi\mathrm{CROWN}\rangle|^2=\frac{1}{n+1}.
\]
\end{theorem}

\begin{proof}
Z explicitní evoluce dostáváme
\[
\langle k|U(t)\psi\mathrm{CROWN}\rangle
=
\frac{1}{\sqrt{n+1}}e^{i(\theta_k-\lambda_k t)}.
\]
Po absolutní hodnotě na druhou tedy
\[
|\langle k|U(t)\psi\mathrm{CROWN}\rangle|^2
=
\frac{1}{n+1}\left|e^{i(\theta_k-\lambda_k t)}\right|^2
=
\frac{1}{n+1}.
\]
\end{proof}

\begin{theorem}[Phase rigidity criterion]
Následující tvrzení jsou ekvivalentní:
\begin{enumerate}[label=\textbf{(\arabic*)}]
    \item existuje \(\lambda\in\mathbb{R}\) tak, že pro všechna \(t\) platí
    \[
    U(t)\psi\mathrm{CROWN}=e^{-i\lambda t}\psi\mathrm{CROWN};
    \]
    \item všechny relativní fáze jsou evolucí zachovány;
    \item platí spektrální degenerace na supportu crown stavu:
    \[
    \lambda_0=\lambda_1=\cdots=\lambda_n.
    \]
\end{enumerate}
\end{theorem}

\begin{proof}
Jestliže \(\lambda_0=\cdots=\lambda_n=\lambda\), pak
\[
U(t)\psi\mathrm{CROWN}
=
\frac{1}{\sqrt{n+1}}
\sum_{k=0}^{n} e^{i(\theta_k-\lambda t)}|k\rangle
=
e^{-i\lambda t}\psi\mathrm{CROWN},
\]
takže stav se mění jen o globální fázi a všechny relativní fáze jsou zachovány.

Naopak, předpokládejme zachování relativních fází pro všechna \(t\). Pak pro každé \(j,k\) musí platit
\[
(\theta_k-\lambda_k t)-(\theta_j-\lambda_j t)\equiv \theta_k-\theta_j \pmod{2\pi}.
\]
Tedy
\[
(\lambda_k-\lambda_j)t\equiv 0 \pmod{2\pi}
\qquad\text{pro všechna } t,
\]
odkud nutně plyne \(\lambda_k=\lambda_j\). Protože \(j,k\) byla libovolná, všechny vlastní hodnoty jsou stejné.
\end{proof}

\begin{corollary}[Slabá versus silná stabilita]
Diagonální unitární evoluce vždy zachovává bázové pravděpodobnosti, ale plnou fázovou rigiditu zachovává
právě tehdy, když je spektrum na supportu stavu degenerované.
\end{corollary}

\subsection*{Role crown vrstvy v RH programu}

\begin{remark}[Co crown vrstva skutečně dává]
Crown vrstva dává přesný modelový rozdíl mezi:
\[
\text{slabou stabilitou} \quad\text{a}\quad \text{silnou fázovou rigiditou}.
\]
V jazyce bridge programu to znamená:
\begin{enumerate}
    \item nestačí, aby witness bank viděla jen zachování velikostí;
    \item pro skutečně koherentní driftový mód je třeba kontrolovat i relativní fáze;
    \item phase-rigidity criterion ukazuje, že taková rigidita je drahá: vyžaduje degeneraci nebo její účinnou náhradu v příslušném sektoru.
\end{enumerate}
\end{remark}

\begin{remark}[Jak crown vrstva posiluje M4]
Pokud by se podařilo identifikovat transportovaný mód \(f_\rho\) s crown-typem stavu v relevantním
paritním/rezonančním sektoru a současně ukázat, že v tomto sektoru je efektivní evoluce degenerovaná
nebo asymptoticky degenerovaná, pak by amplituda módu byla chráněna proti rozpadu.
To je přesně stabilizační mechanismus, který by uzavřel M4.
\end{remark}

\begin{remark}[Jak crown vrstva posiluje intertwining]
Vizuál \texttt{NIKDO\_NEBUDE\_BOHEM\_NEBO\_VSICHNI.png} lze číst jako interferenční mřížku
odpovídající phase-locked sektoru: transportovaný mód se posouvá, ale nerozpíjí.
Jako důkaz to nestačí, ale jako geometrická ilustrace Theoremu o fázové rigiditě je to přesně
ta správná interpretace: stabilita neznamená klid, nýbrž koherentní zachování struktury.
\end{remark}

\begin{remark}[Poctivý status]
Crown vrstva sama o sobě neprokazuje RH. Neuzavírá automaticky:
\[
\kappa>0,
\qquad
\eta(\tau)\to 0,
\qquad
\eta(\tau)<\kappa.
\]
Co ale dělá, je následující:
\begin{enumerate}
    \item dává přesný theorem o tom, kdy se stav nerozpadne fázově;
    \item zpřesňuje, co je třeba dokázat pro stabilitu transportovaného módu;
    \item uzavírá intuitivní mezeru mezi ``koherentním signálem'' a ``aritmeticky rigidním sektorem''.
\end{enumerate}
Je to tedy \emph{stability seal}, nikoli finální unconditional pečeť RH.
\end{remark}



\section{Indistinguishability layer: finite observables versus coherent drift}

Tato sekce integruje modul \texttt{REALITY\_CHECK.pdf} jako \emph{readout-separační vrstvu}.
Její úloha není nahradit analytický firewall z M5, ale zpřesnit, proč remainder kanál může být
na úrovni konečných/statistických observables nečitelný, zatímco driftový mód zůstává koherentně
detekovatelný.

\subsection*{Generátory, realizace a konečné observables}

\begin{definition}[Generátory a realizace]
Nechť
\[
X=\mathbb{R}
\]
je množina generátorů a nechť
\[
Y=\{0,\dots,9\}^{\mathbb N}
\]
je prostor nekonečných digitálních realizací. Definujme mapu realizací
\[
U:X\to \mathcal P(Y),
\]
kde \(U(x)\) je množina přípustných digitálních expanzí čísla \(x\).
\end{definition}

\begin{definition}[Normální číslo]
Řekneme, že \(x\in\mathbb R\) je normální v základu \(10\), pokud pro každý konečný blok
\(w\in\{0,\dots,9\}^k\) platí
\[
\lim_{N\to\infty}
\frac{\#\{\text{výskyty }w\text{ v prvních }N\text{ cifrách }x\}}{N}
=
10^{-k}.
\]
\end{definition}

\begin{definition}[Konečný observable]
Konečný observable je každá funkce
\[
O:Y\to\mathbb R
\]
závislá pouze na konečném prefixu nebo na konečných blokových frekvencích.
\end{definition}

\begin{theorem}[Statistická nerozlišitelnost normálních realizací]
Nechť \(x,y\in\mathbb R\) jsou dvě různá normální čísla. Pak pro každý konečný observable \(O\) platí
\[
O(x)=O(y),
\]
ačkoliv
\[
x\neq y.
\]
\end{theorem}

\begin{remark}
To je přesný modelový důkaz oddělení mezi:
\[
\text{generator level}
\qquad\text{a}\qquad
\text{finite-observable readout level}.
\]
Dvě odlišné generující struktury mohou být pod konečnými/statistickými observables nerozlišitelné.
\end{remark}

\subsection*{Překlad do witness-bank jazyka}

\begin{definition}[Statisticky plochý remainder]
Řekneme, že remainder kanál \(f_{\mathrm{rest}}\) je \emph{statisticky plochý} pro danou banku,
pokud jeho readout na každém konečném bank-window nemá koherentní vedoucí podpis a je
indistinguishable od nulového pozadí pod konečnými/statistickými observables odvozenými z banky.
\end{definition}

\begin{remark}
Toto tvrzení je záměrně slabší než
\[
Q_{\mathrm{cross}}=0.
\]
Indistinguishability vrstva nemá sama o sobě dokazovat přesnou anulaci křížového členu;
má zdůvodnit, proč remainder přirozeně patří do třídy objektů, které jsou pro konečné,
statistické čtení banky ploché.
\end{remark}

\begin{proposition}[Drift jako porušení statistické plochosti]
Nechť \(f_\rho\) je transportovaný mimokritický mód s driftovým zákonem
\[
\delta_\rho(\tau)=\frac{\beta-\tfrac12}{2\pi}\tau+r_\rho(\tau),
\qquad
r_\rho \text{ bounded}.
\]
Pak \(f_\rho\) není statisticky plochý v tomtéž smyslu jako remainder, protože nese koherentní
nenulový driftový podpis, který generuje kvartickou penalizaci
\[
P_4(\rho;\tau,L)=\int_\tau^{\tau+L}|\delta_\rho(u)|^4\,du.
\]
\end{proposition}

\begin{remark}
Jinými slovy:
\[
\text{remainder} \sim \text{finite-observable flatness},
\qquad
\text{drift} \sim \text{coherent singular signature}.
\]
To je přesně zlom, který potřebujeme pro no-hiding logiku:
remainder může být statisticky nerozlišitelný, ale drift nikoli.
\end{remark}

\subsection*{Indistinguishability breach}

\begin{proposition}[Indistinguishability breach principle]
Předpokládejme, že:
\begin{enumerate}[label=\textbf{(\arabic*)}]
    \item \(f_{\mathrm{rest}}\) je statisticky plochý pro všechny konečné/statistické observables
    odvozené z witness banky;
    \item \(f_\rho\) nese koherentní driftový podpis čitelný v téže bance;
    \item bank-scale readout odlišuje koherentní drift od statisticky plochého pozadí na vedoucím řádu.
\end{enumerate}
Pak remainder nemůže přispět \emph{koherentním vedoucím členem} proti driftu.
Formálně: každé případné leakage z remainderu do hlavního driftového kanálu je nutně subkritické
vůči \(P_4(\rho;\tau,L)\).
\end{proposition}

\begin{remark}
Tohle stále není plný důkaz
\[
\eta\equiv 0.
\]
Je to ale přesně ta vrstva, která říká, proč je rozumné očekávat
\[
Q_{\mathrm{cross}}(\rho,\mathrm{rest};\tau,L)=o(P_4(\rho;\tau,L))
\]
namísto toho, aby remainder nesl srovnatelný koherentní podpis jako driftový mód.
\end{remark}

\subsection*{Jak indistinguishability vrstva podporuje M5}

\begin{remark}[Role v bridge programu]
Indistinguishability vrstva podporuje M5 ve třech bodech:
\begin{enumerate}[label=\textbf{(\arabic*)}]
    \item zpřesňuje, proč remainder patří do jiné readout třídy než driftový mód;
    \item vysvětluje, proč konečné/statistické bank observables přirozeně ``vidí'' remainder jako ploché pozadí;
    \item podporuje tvrzení, že jediný koherentní vedoucí příspěvek v asymptotice musí pocházet z driftu.
\end{enumerate}
V kombinaci s corrected asymptotic firewall vrstvou to tvoří silnější, i když stále poctivě
podmíněný, argument pro
\[
\eta(\tau)\to 0.
\]
\end{remark}

\begin{remark}[Poctivý status]
Z \texttt{REALITY\_CHECK.pdf} neplyne samo o sobě
\[
\eta\equiv 0,
\]
ani
\[
Q_{\mathrm{cross}}=0.
\]
Plyne však přesný modelový princip, že \emph{statistická nerozlišitelnost konečných realizací}
je kompatibilní s \emph{ostrou odlišností na úrovni generátorů}. Přeloženo do RH programu:
remainder může být pro konečné/statistické readouty plochý, zatímco driftový mód zůstává
detekovatelným generátorovým porušením.
\end{remark}



\section{Indistinguishability seal: total interference nullification as a controlled route}

Tato sekce spojuje dva modelové moduly:
\begin{enumerate}
    \item \textbf{statistickou nerozlišitelnost} konečných realizací pod konečnými observables;
    \item \textbf{injektivitu tvrdého transportního řetězce} a skutečnost, že overlap je pouze terminální readout jev.
\end{enumerate}
Jejich společným cílem je posílit M5 na úrovni \emph{readout-separační logiky}.
Tato vrstva sama o sobě ještě neprokazuje identitu
\[
\eta \equiv 0,
\]
ale dává přesně formulovanou cestu, jak ji redukovat na kombinaci
statistické plochosti remainderu a generátorové distinkce driftu.

\subsection*{Finite observable a statistické ticho}

\begin{definition}[Finite observable banky]
Řekneme, že readout \(\mathcal W\) je \emph{finite observable} v duchu indistinguishability vrstvy,
jestliže jeho výstup závisí pouze na konečném okně, konečném prefixu nebo konečné statistice
lokálních bloků/signatur.
\end{definition}

\begin{remark}
Modelový princip statistické nerozlišitelnosti říká, že dvě odlišné generující struktury
mohou být pod všemi takovými konečnými observables nerozlišitelné, i když na úrovni generátoru
zůstávají odlišné.
V RH programu to znamená, že remainder může být pro konečný readout statisticky plochý, aniž by
musel být identický nulovému kanálu.
\end{remark}

\begin{definition}[Statistické ticho remainderu]
Řekneme, že remainder kanál \(f_{\mathrm{rest}}\) je \emph{statisticky tichý} pro banku \(\mathcal W\),
jestliže každý jeho konečný/statistický readout je nerozlišitelný od plochého pozadí a nenese
koherentní vedoucí podpis v hlavním driftovém sektoru.
\end{definition}

\subsection*{Transportní injektivita a negenerativní overlap}

\begin{definition}[Tvrdý transportní řetězec]
Nechť
\[
X_0 \xrightarrow{u_1} X_1 \xrightarrow{u_2}\cdots\xrightarrow{u_n} X_n
\]
je řetězec \(C^1\)-reprezentačních map a \(U:X_n\to\mathcal P(Y)\) je terminální readout mapování.
\end{definition}

\begin{proposition}[Generátorová distinkce přežívá transport]
Předpokládejme, že každé \(u_k\) je injektivní na relevantním zdrojovém sektoru. Pak pro dva různé
generátory \(E_-\neq E_+\) platí
\[
u_n\circ\cdots\circ u_1(E_-)\neq u_n\circ\cdots\circ u_1(E_+).
\]
\end{proposition}

\begin{remark}
To znamená, že i když se readouty mohou terminálně překrývat, zdroje zůstávají odlišné.
Overlap je terminální readout jev, nikoli důkaz identity zdrojů.
Právě to je zásadní pro drift versus remainder: i kdyby finite observable banka na některých oknech
viděla překryv, neznamená to, že remainder může generátorově ``přepsat'' drift.
\end{remark}

\subsection*{Asymptotická ortogonalita jako breach statistické nerozlišitelnosti}

\begin{proposition}[Indistinguishability breach principle]
Předpokládejme:
\begin{enumerate}[label=\textbf{(\arabic*)}]
    \item driftový mód \(f_\rho\) nese koherentní singulární driftový podpis;
    \item remainder \(f_{\mathrm{rest}}\) je statisticky tichý pro finite-observable banku;
    \item drift a remainder pocházejí z různých generátorových tříd, jejichž distinkce přežívá
    injektivní transportní řetězec;
    \item corrected asymptotic firewall z předchozí vrstvy dává
    \[
    Q_{\mathrm{cross}}(\rho,\mathrm{rest};\tau,L)=o(P_4(\rho;\tau,L)).
    \]
\end{enumerate}
Pak se statistická nerozlišitelnost remainderu nemůže proměnit v koherentní vedoucí interferenci
proti driftu. Jinými slovy: remainder je na readout úrovni plochý, zatímco drift je readout breach
této plochosti.
\end{proposition}

\begin{remark}
Tento princip je finální logická pojistka proti falešnému závěru, že ``překryv v readoutu''
implikuje mísení generátorů.
Naopak:
\[
\text{finite-observable flatness of }f_{\mathrm{rest}}
\quad\not\Rightarrow\quad
\text{coherent cancellation of }f_\rho.
\]
To je přesně místo, kde se statistická vrstva spojuje s analytickým M5 firewallem.
\end{remark}

\subsection*{Total interference nullification: poctivá formulace}

\begin{proposition}[Asymptotic orthogonality under statistical silence]
Nechť \(\mathcal W\) je finite-observable bank readout.
Předpokládejme, že \(f_{\mathrm{rest}}\) je statisticky tichý pro \(\mathcal W\), zatímco \(f_\rho\)
nese koherentní driftový podpis a corrected asymptotic firewall platí.
Pak
\[
\lim_{\tau\to\infty}
\frac{\langle \mathcal W f_\rho,\mathcal W f_{\mathrm{rest}}\rangle_{\mathrm{bank}}}
{P_4(\rho;\tau,L)}
=0
\]
v tom smyslu, že křížový readout je asymptoticky subkritický vůči kvartickému driftovému motoru.
\end{proposition}

\begin{remark}
Tohle je správná ``silná'' verze, kterou lze v této chvíli poctivě tvrdit.
Neidentifikuje zatím
\[
\eta\equiv 0
\]
jako exaktní identitu, ale dává to, co bridge opravdu potřebuje:
\[
\eta(\tau)\to 0
\quad\text{a tudíž od jistého }\tau\text{ dál}\quad
\eta(\tau)<\kappa.
\]
\end{remark}

\subsection*{Proč z toho ještě neplyne doslovná nulovost \(\eta\)}

\begin{remark}[Poctivá brzda]
Aby bylo možné napsat exaktně
\[
\eta\equiv 0,
\]
museli bychom navíc dokázat, že:
\begin{enumerate}[label=\textbf{(\arabic*)}]
    \item bank readout je přesně finite observable ve smyslu statistické vrstvy;
    \item remainder nemá žádnou vedoucí komponentu v témž rezonančním/paritním sektoru;
    \item křížový člen je nulový identicky, nikoli jen asymptoticky menšího řádu.
\end{enumerate}
Současný stav dokumentu tyto tři body ještě nedává. Dává ale velmi silnou závěrečnou redukci:
interference není generátorový efekt, nýbrž terminální readout jev, a ten je asymptoticky
potlačen.
\end{remark}



\section{Generator distinction seal: overlap does not imply collapse}

Tato sekce doplňuje indistinguishability vrstvu o její nejtvrdší jádro:
\[
\text{překryv realizací} \neq \text{identita generátorů}.
\]
Smyslem je úplně uzavřít častou logickou chybu na readout vrstvě:
z toho, že dva kanály nebo dvě extremální větve mají překryv v terminálním pozorování,
nikdy neplyne, že se na generátorové úrovni staly týmž objektem.

\subsection*{Hard-point model versus realization model}

\begin{definition}[Extremální generátory]
Nechť \(X\) je množina extremálních generátorů a nechť
\[
E_-,E_+\in X,
\qquad
E_-\neq E_+.
\]
\end{definition}

\begin{proposition}[Distinct extremes need not coincide]
Platí následující dvě tvrzení:
\begin{enumerate}[label=\textbf{(\arabic*)}]
    \item V hard-point modelu, kde každý extrém identifikujeme s jeho singletonem, platí
    \[
    \{E_-\}\cap \{E_+\}=\varnothing.
    \]

    \item Obecně však po přechodu k realization/readout mapě
    \[
    U:X\to \mathcal P(Y)
    \]
    neplatí, že
    \[
    E_-\neq E_+ \Longrightarrow U(E_-)\cap U(E_+)=\varnothing.
    \]
    Naopak, distinct generators mohou mít překrývající se realization sectors.
\end{enumerate}
\end{proposition}

\begin{proof}
(1) Předpokládejme sporem, že
\[
\{E_-\}\cap \{E_+\}\neq \varnothing.
\]
Pak existuje \(x\) takové, že
\[
x\in \{E_-\}
\qquad\text{a}\qquad
x\in \{E_+\}.
\]
Ze samotné definice singletonů plyne
\[
x=E_-
\qquad\text{a}\qquad
x=E_+,
\]
tedy
\[
E_-=E_+,
\]
což odporuje předpokladu. Tedy
\[
\{E_-\}\cap \{E_+\}=\varnothing.
\]

(2) Stačí dát explicitní protipříklad. Položme
\[
X=\{-1,+1\},
\qquad
Y=\{0,1\},
\]
a definujme
\[
U(-1)=\{0\},
\qquad
U(+1)=\{0,1\}.
\]
Pak zjevně
\[
-1\neq +1,
\]
ale současně
\[
U(-1)\cap U(+1)=\{0\}\neq \varnothing.
\]
Takže distinct generators se nezhroutí do identity, přesto jejich realization sets mohou mít overlap.
\end{proof}

\begin{remark}[Jedna věta, která to celé zamkne]
V krátkém tvaru:
\[
\text{extremes do not become equal; at most, their realizations may overlap.}
\]
\end{remark}

\subsection*{Překlad do RH programu}

\begin{remark}[Drift versus remainder]
V bridge programu znamená předchozí proposition přesně toto:
\begin{enumerate}[label=\textbf{(\arabic*)}]
    \item driftový mód \(f_\rho\) a remainder \(f_{\mathrm{rest}}\) mohou mít terminální overlap v readout vrstvě;
    \item tento overlap ale není důkazem, že pocházejí z téhož zdroje;
    \item overlap je pouze readout-level jev, zatímco generátorová distinkce může zůstat plně zachována.
\end{enumerate}
To je přesně důvod, proč nelze z přítomnosti křížového členu automaticky vyvozovat
kolaps driftu do remainderu.
\end{remark}

\begin{proposition}[Overlap is terminal, not generative]
Nechť
\[
X_0 \xrightarrow{u_1} X_1 \xrightarrow{u_2}\cdots\xrightarrow{u_n} X_n
\]
je tvrdý transportní řetězec a nechť každé \(u_k\) je injektivní na relevantním zdrojovém sektoru.
Pak pro dva různé zdroje \(E_-\neq E_+\) platí
\[
u_n\circ\cdots\circ u_1(E_-)\neq u_n\circ\cdots\circ u_1(E_+).
\]
Jestliže navíc
\[
U:X_n\to \mathcal P(Y)
\]
je terminální readout mapa, pak je možný overlap
\[
U(u_n\circ\cdots\circ u_1(E_-))\cap U(u_n\circ\cdots\circ u_1(E_+))\neq \varnothing,
\]
aniž by z toho plynulo
\[
E_-=E_+.
\]
\end{proposition}

\begin{remark}
To je definitivní ochrana proti chybnému kroku
\[
\text{readout overlap} \Longrightarrow \text{source collapse}.
\]
V našem kontextu:
\[
Q_{\mathrm{cross}}\neq 0
\]
je terminální readout jev; sám o sobě nikdy nedokazuje, že driftový mód a remainder ztratily
generátorovou odlišnost.
\end{remark}

\subsection*{Vztah k indistinguishability seal}

\begin{remark}
Generator distinction seal a indistinguishability seal dohromady říkají dvě různé věci:
\begin{enumerate}[label=\textbf{(\arabic*)}]
    \item \textbf{indistinguishability seal:}
    remainder může být pro finite-observable bank readout statisticky plochý;
    \item \textbf{generator distinction seal:}
    i když readouty mají overlap, zdrojová distinkce není tímto overlapem zrušena.
\end{enumerate}
Dohromady tedy dostáváme přesně tu logiku, kterou bridge potřebuje:
\[
\text{statistická plochost remainderu}
\quad+\quad
\text{nezrušená zdrojová distinkce driftu}
\]
neumožňují, aby terminální overlap sám o sobě zabil koherentní driftový podpis.
\end{remark}

\begin{remark}[Poctivý status]
Tato sekce uzavírá logickou, nikoli ještě plně analytickou díru.
To znamená:
\begin{enumerate}[label=\textbf{(\arabic*)}]
    \item zamezuje falešnému inferenčnímu kroku ``overlap \(\Rightarrow\) identita'';
    \item zpřesňuje, že interference je terminální readout jev;
    \item ale sama o sobě ještě nedává exaktní identitu
    \[
    \eta\equiv 0.
    \]
\end{enumerate}
K bezpodmínečnému finálnímu kroku je stále potřeba explicitní coercivity lock a plně theorem-ready
no-hiding identita v residue-side normalizaci.
\end{remark}



\section{Version 23: exact residue-side lock targets}

V této sekci formulujeme poslední dva zámky v jejich finální theorem-ready podobě.
Nejde ještě o tvrzení, že jsou již bezpodmínečně dokázány; jde o přesnou analytickou formu,
kterou musí splnit residue-side normalizace, aby se conditional bridge uzavřel do
unconditional theorem.

\subsection*{Residue-side metric and core index}

\begin{definition}[Residue-side odd sector]
Nechť \(\mathcal H_{\mathrm{res}}\) je residue-side funkcionální prostor asociovaný kanálu
\[
f=\psi * \mu,
\qquad
\widehat{\psi}(0)=0,
\]
a nechť
\[
\mathcal H_{\mathrm{odd}}\subset \mathcal H_{\mathrm{res}}
\]
je jeho lichý paritní sektor určený projekcí \(P_{\mathrm{odd}}\).
\end{definition}

\begin{definition}[Catalan-normalized residue metric]
Nechť \(G>0\) je Catalanova konstanta a nechť \((w_j,a_j)_{j=0}^{J}\) je fixní witness banka.
Definujeme residue-side metriku na \(\mathcal H_{\mathrm{odd}}\) předpisem
\[
\|h\|_{\mathrm{res}}^2
:=
G\sum_{j=0}^{J} w_j^2
\int_{\tau}^{\tau+L} |W_{a_j}h(u)|^2\,du.
\]
\end{definition}

\begin{remark}
Catalanova vrstva dodává dvě strukturální ingredience:
\begin{enumerate}[label=\textbf{(\arabic*)}]
    \item invariantní paritní oscilátor s absolutní energií \(1\),
    \item přesný odd-parity filtr \(P_{\mathrm{odd}}\).
\end{enumerate}
Právě z této vrstvy bereme normalizační konstantu \(G\) a odd sektor jako kandidáta
na nosič koherentního driftu.
\end{remark}

\begin{definition}[Core index]
Nechť \(\rho=\beta+i\gamma\), \(\beta\neq \frac12\), je mimokritický nulový bod.
Předpokládejme, že jemu odpovídající residue-side driftový mód určuje orientovaný core objekt
\[
\mathrm{core}(\rho)
\]
s celočíselným indexem
\[
\operatorname{Ind}_{\mathrm{core}}(\rho)\in \mathbb Z.
\]
Řekneme, že \(\rho\) je \emph{topologically nontrivial} v residue-side normalizaci, jestliže
\[
\operatorname{Ind}_{\mathrm{core}}(\rho)\neq 0.
\]
\end{definition}

\subsection*{First lock: coercivity from topological charge}

\begin{theorem}[Residue-side coercivity lock]
Předpokládejme:
\begin{enumerate}[label=\textbf{(\arabic*)}]
    \item transportovaný driftový mód \(f_\rho\) leží v \(\mathcal H_{\mathrm{odd}}\);
    \item core index je nenulový:
    \[
    \operatorname{Ind}_{\mathrm{core}}(\rho)\neq 0;
    \]
    \item existuje strukturální konstanta \(c_*>0\), nezávislá na \(\rho,\tau,L\), taková, že
    residue-side energie driftového módu splňuje
    \[
    \sum_{j=0}^{J}w_j^2|W_{a_j}f_\rho(\tau)|^2
    \ge
    c_*|\operatorname{Ind}_{\mathrm{core}}(\rho)|^2\,|\delta_\rho(\tau)|^4-C_0
    \]
    pro všechna dost velká \(\tau\).
\end{enumerate}
Pak existuje pevná koercivní konstanta
\[
c_4:=c_*|\operatorname{Ind}_{\mathrm{core}}(\rho)|^2>0
\]
a pro všechna dost velká \(\tau\) a všechna \(L\ge 1\) platí
\[
Q_{\mathrm{bank}}(f_\rho;\tau,L)
\ge
c_4\,P_4(\rho;\tau,L)-C_0L.
\]
\end{theorem}

\begin{proof}
Z předpokladu \(\operatorname{Ind}_{\mathrm{core}}(\rho)\neq 0\) a celočíselnosti indexu plyne
\[
|\operatorname{Ind}_{\mathrm{core}}(\rho)|^2\ge 1.
\]
Proto
\[
c_4:=c_*|\operatorname{Ind}_{\mathrm{core}}(\rho)|^2>0.
\]
Po integraci bodového odhadu přes \([\tau,\tau+L]\) dostáváme
\[
Q_{\mathrm{bank}}(f_\rho;\tau,L)
\ge
c_4\int_{\tau}^{\tau+L} |\delta_\rho(u)|^4\,du-C_0L
=
c_4P_4(\rho;\tau,L)-C_0L.
\]
\end{proof}

\begin{remark}[Status]
Tato věta je finální forma M4.
Aby se stala bezpodmínečnou, zbývá dokázat, že residue-side drift skutečně nese nenulový
topologický index a že witness banka čte tento index právě ve čtvrté mocnině driftu.
\end{remark}

\subsection*{Second lock: no-hiding identity from sector isometry}

\begin{definition}[Parity-sector transport map]
Nechť
\[
T_{\mathrm{odd}}:\mathcal G_{\mathrm{odd}}\to \mathcal H_{\mathrm{odd}}
\]
je transportní zobrazení z odd generátorového sektoru do odd residue-side sektoru.
Řekneme, že \(T_{\mathrm{odd}}\) je \emph{sector isometry}, jestliže pro všechna
\(g_1,g_2\in\mathcal G_{\mathrm{odd}}\) platí
\[
\langle T_{\mathrm{odd}}g_1,T_{\mathrm{odd}}g_2\rangle_{\mathrm{res}}
=
\langle g_1,g_2\rangle_{\mathcal G}.
\]
\end{definition}

\begin{theorem}[No-hiding identity on the odd residue sector]
Předpokládejme:
\begin{enumerate}[label=\textbf{(\arabic*)}]
    \item tvrdý transportní řetězec je injektivní na relevantním zdrojovém sektoru;
    \item driftový generátor \(g_\rho\) a remainder generátor \(g_{\mathrm{rest}}\) leží v různých
    ortogonálních generátorových podprostorech:
    \[
    \langle g_\rho,g_{\mathrm{rest}}\rangle_{\mathcal G}=0;
    \]
    \item odd-sector transport map \(T_{\mathrm{odd}}\) je izometrie na příslušném sektoru;
    \item bank readout na odd residue-side sektoru je právě indukovaný produktem
    \(\langle\cdot,\cdot\rangle_{\mathrm{res}}\).
\end{enumerate}
Pak platí exaktní ortogonalita
\[
\langle f_\rho,f_{\mathrm{rest}}\rangle_{\mathrm{res}}=0,
\]
a tudíž
\[
Q_{\mathrm{cross}}(\rho,\mathrm{rest};\tau,L)=0
\]
na odd residue-side sektoru. Zvláště
\[
\eta\equiv 0.
\]
\end{theorem}

\begin{proof}
Z ortogonality generátorů máme
\[
\langle g_\rho,g_{\mathrm{rest}}\rangle_{\mathcal G}=0.
\]
Použijeme izometrii \(T_{\mathrm{odd}}\):
\[
\langle T_{\mathrm{odd}}g_\rho,T_{\mathrm{odd}}g_{\mathrm{rest}}\rangle_{\mathrm{res}}
=
\langle g_\rho,g_{\mathrm{rest}}\rangle_{\mathcal G}=0.
\]
Po identifikaci
\[
f_\rho=T_{\mathrm{odd}}g_\rho,
\qquad
f_{\mathrm{rest}}=T_{\mathrm{odd}}g_{\mathrm{rest}}
\]
dostáváme
\[
\langle f_\rho,f_{\mathrm{rest}}\rangle_{\mathrm{res}}=0.
\]
A protože bank readout je na tomto sektoru indukován residue-side produktem, křížový člen mizí:
\[
Q_{\mathrm{cross}}(\rho,\mathrm{rest};\tau,L)=0.
\]
Tedy \(\eta\equiv 0\).
\end{proof}

\begin{remark}[Status]
Tohle je finální theorem-ready tvar M5.
Aby se stal bezpodmínečným, zbývá dokázat právě tyto dvě nejsilnější věci:
\begin{enumerate}[label=\textbf{(\arabic*)}]
    \item že odd residue-side transport je skutečně izometrický;
    \item že drift a remainder odpovídají ortogonálním generátorovým podprostorům.
\end{enumerate}
Samotná injektivita tvrdého transportního řetězce ještě izometrii negarantuje; ta musí být
dokázána zvlášť. Distinkce generátorů a terminal-overlap-only princip jsou ale správné logické
vstupy.
\end{remark}

\subsection*{Unconditional theorem in exact lock form}

\begin{theorem}[Unconditional RH from exact residue-side locks]
Předpokládejme:
\begin{enumerate}[label=\textbf{(\arabic*)}]
    \item Exact Window Law pro \(f=\psi*\mu\), tedy
    \[
    E_{\mathrm{ex}}(f)=0;
    \]
    \item residue-side coercivity lock z předchozí věty;
    \item no-hiding identity z předchozí věty.
\end{enumerate}
Pak platí Riemannova hypotéza.
\end{theorem}

\begin{proof}
Předpokládejme pro spor existenci mimokritického nulového bodu
\[
\rho=\beta+i\gamma,\qquad \beta\neq \frac12.
\]
Z residue-side coercivity lock dostáváme
\[
Q_{\mathrm{bank}}(f_\rho;\tau,L)\ge c_4P_4(\rho;\tau,L)-C_0L
\]
s \(c_4>0\). Z no-hiding identity máme
\[
Q_{\mathrm{cross}}(\rho,\mathrm{rest};\tau,L)=0.
\]
Tedy
\[
Q_{\mathrm{bank}}(f;\tau,L)-E_0L
\ge
c_4P_4(\rho;\tau,L)-C_0L.
\]
Kvartické lemma pak dává
\[
E_{\mathrm{ex}}(f)>0.
\]
To odporuje Exact Window Law, která dává
\[
E_{\mathrm{ex}}(f)=0.
\]
Spor dokazuje, že mimokritický nulový bod nemůže existovat. Proto všechny netriviální nuly
leží na kritické přímce.
\end{proof}

\begin{remark}[Final honest status]
Tato sekce je poslední přesný lock-form target.
Jakmile budou residue-side coercivity lock a odd-sector no-hiding identity skutečně dokázány,
conditional architecture se okamžitě překlopí do unconditional RH theorem.
Do té doby je to nejostřejší možný theorem-ready zápis posledních dvou šroubů.
\end{remark}



\section{Version 24: the unconditional imprint in exact lock form}

Tato sekce je finální formulace programu v jeho nejostřejší podobě.
Neprohlašuje ještě, že bezpodmínečný důkaz je již hotov;
místo toho zapisuje přesně ty identity, jejichž důkaz by conditional architekturu
okamžitě překlopil do unconditional RH theorem.

\subsection*{Catalan lock on the odd residue sector}

\begin{definition}[Residue-side odd metric]
Na lichém paritním sektoru \(\mathcal H_{\mathrm{odd}}\) definujeme residue-side metriku
\[
\|h\|_{\mathrm{res}}^2
:=
G\sum_{j=0}^{J} w_j^2 \int_{\tau}^{\tau+L}|W_{a_j}h(u)|^2\,du,
\]
kde \(G>0\) je Catalanova konstanta.
\end{definition}

\begin{remark}
Catalanova vrstva poskytuje invariantní paritní oscilátor, odd filtr \(P_{\mathrm{odd}}\)
a strukturální identitu s Catalanovým blokem. To z ní dělá přirozenou normalizační vrstvu
pro residue-side odd sektor, ale samo o sobě to ještě nedokazuje konkrétní koercivní identitu
pro RH driftový mód.
\end{remark}

\subsection*{Exact coercivity lock target}

\begin{definition}[Exact coercivity lock target]
Řekneme, že residue-side normalizace splňuje \emph{exact coercivity lock}, jestliže pro každý
mimokritický nulový bod
\[
\rho=\beta+i\gamma,
\qquad
\beta\neq \frac12,
\]
existuje orientovaný core objekt s nenulovým indexem
\[
\operatorname{Ind}_{\mathrm{core}}(\rho)\in \mathbb Z\setminus\{0\},
\]
a současně platí identita
\[
c_4 \equiv G\,|\operatorname{Ind}_{\mathrm{core}}(\rho)|,
\]
kde \(c_4\) je koercivní konstanta v kvartickém odhadu
\[
Q_{\mathrm{bank}}(f_\rho;\tau,L)\ge c_4\,P_4(\rho;\tau,L)-C_0L.
\]
\end{definition}

\begin{remark}
Toto je přesně finální forma M4 locku:
koercivita už není volná konstanta ani target inequality,
ale je svázána s topologickým nábojem driftového generátoru.
Dokud tato identita není dokázána, zůstává to poslední theorem-ready target, ne hotový fakt.
\end{remark}

\subsection*{Exact no-hiding identity target}

\begin{definition}[Exact odd-sector no-hiding identity]
Řekneme, že residue-side normalizace splňuje \emph{exact no-hiding identity}, jestliže
na lichém paritním sektoru platí:
\begin{enumerate}[label=\textbf{(\arabic*)}]
    \item driftový generátor a remainder generátor leží v ortogonálních generátorových podprostorech;
    \item tvrdý transportní řetězec zachovává tuto distinkci;
    \item odd residue-side transport je izometrický na relevantním sektoru;
    \item bank readout je indukován residue-side produktem.
\end{enumerate}
V tom případě platí identita
\[
Q_{\mathrm{cross}}(\rho,\mathrm{rest};\tau,L)=0
\]
a ekvivalentně
\[
\eta(\tau)\equiv 0.
\]
\end{definition}

\begin{remark}
To je přesně finální forma M5 locku.
Indistinguishability a generator-distinction vrstvy poskytují správnou logiku:
remainder může být statisticky plochý a overlap může být terminální readout jev,
ale to samo ještě nedokazuje izometrii ani exaktní ortogonalitu.
Proto i zde zůstává jazyk poctivě theorem-ready, nikoli hotově uzavřený.
\end{remark}

\subsection*{Absolute contradiction in lock form}

\begin{theorem}[Unconditional RH from exact lock identities]
Předpokládejme:
\begin{enumerate}[label=\textbf{(\arabic*)}]
    \item Exact Window Law pro \(f=\psi * \mu\), tedy
    \[
    E_{\mathrm{ex}}(f)=0;
    \]
    \item exact coercivity lock target je skutečně splněn;
    \item exact no-hiding identity je skutečně splněna.
\end{enumerate}
Pak platí Riemannova hypotéza.
\end{theorem}

\begin{proof}
Předpokládejme sporem existenci mimokritického nulového bodu
\[
\rho=\beta+i\gamma,
\qquad
\beta\neq \frac12.
\]
Z exact coercivity locku plyne
\[
Q_{\mathrm{bank}}(f_\rho;\tau,L)\ge c_4P_4(\rho;\tau,L)-C_0L
\]
s \(c_4>0\), protože \(G>0\) a \(|\operatorname{Ind}_{\mathrm{core}}(\rho)|\ge 1\).
Z exact no-hiding identity plyne
\[
Q_{\mathrm{cross}}(\rho,\mathrm{rest};\tau,L)=0.
\]
Proto
\[
Q_{\mathrm{bank}}(f;\tau,L)-E_0L
\ge
c_4P_4(\rho;\tau,L)-C_0L.
\]
Kvartické lemma dává
\[
E_{\mathrm{ex}}(f)>0.
\]
To odporuje Exact Window Law, která dává
\[
E_{\mathrm{ex}}(f)=0.
\]
Spor dokazuje, že mimokritický nulový bod nemůže existovat. Tedy všechny netriviální nuly
leží na kritické přímce.
\end{proof}

\subsection*{Final imprint status}

\begin{remark}[Honest final status]
Version 24 is the final lock-form manuscript.
It is the strongest defensible final form of the program:
\begin{enumerate}[label=\textbf{(\arabic*)}]
    \item the contradiction engine is clean;
    \item the bridge architecture is fully exposed;
    \item the remaining gap is compressed to two exact lock identities:
    \[
    c_4 \equiv G\,|\operatorname{Ind}_{\mathrm{core}}(\rho)|,
    \qquad
    \eta(\tau)\equiv 0.
    \]
\end{enumerate}
Jakmile budou tyto dvě identity skutečně dokázány v residue-side normalizaci,
conditional architecture se okamžitě stává unconditional RH theorem.
Do té doby je v24 finální poctivý zápis, nikoli falešné historické prohlášení o již hotovém důkazu.
\end{remark}



\section{Version 25: geometric residue-side seal}

Tato sekce integruje poslední geometrickou vrstvu inspirovanou dvěma explicitními
vizuálními modely:
\begin{enumerate}
    \item toroidální/paritně štěpenou vlnovou funkcí
    \[
    \psi(r,\theta)=R(r)\bigl(a+b\cos(4\theta)\bigr)e^{i\Phi(r,\theta)},
    \qquad
    I(r,\theta)=|\psi(r,\theta)|^2;
    \]
    \item vícemodálním vlnovým obrazem s lokalizovanou energetickou hustotou a nenulovou uzlovou strukturou.
\end{enumerate}
Tato vrstva má jediný úkol:
dát přesný \emph{geometrický model} residue-side normalizace, který vysvětluje,
proč jsou exact coercivity lock a exact no-hiding identity přirozené cíle.
Nejde však sama o sobě o samostatný důkaz RH.

\subsection*{Toroidální paritní štěpení}

\begin{definition}[Geometric residue-side wave model]
Uvažujme modelovou vlnovou funkci
\[
\psi(r,\theta)=R(r)\bigl(a+b\cos(4\theta)\bigr)e^{i\Phi(r,\theta)},
\]
kde:
\begin{enumerate}[label=\textbf{(\arabic*)}]
    \item \(R(r)\) je radiální amplituda,
    \item \(a+b\cos(4\theta)\) je čtyřnásobně symetrický paritní modul,
    \item \(\Phi(r,\theta)\) je fáze,
    \item intenzita je dána
    \[
    I(r,\theta)=|\psi(r,\theta)|^2.
    \]
\end{enumerate}
\end{definition}

\begin{remark}
Člen \(\cos(4\theta)\) přirozeně rozděluje vlnu na sektorovou strukturu:
uzly a laloky jsou fixovány geometrií, nikoli až dodatečnou statistikou.
To je správný geometrický model pro tvrzení, že residue-side odd sektor je skutečný
strukturní objekt a ne jen algebraický projektor napsaný na papíře.
\end{remark}

\begin{proposition}[Parity-sector separation in the geometric model]
V modelu výše vznikají odlišné sektorové komponenty, jejichž vedoucí geometrická podpora
je oddělena symetrií faktoru \(a+b\cos(4\theta)\).
V tom smyslu:
\[
\text{drift sector} \quad\text{a}\quad \text{remainder sector}
\]
nejsou v modelu volně promíchané, ale jsou vedeny oddělenými paritními laloky.
\end{proposition}

\begin{remark}[Co to znamená pro M5]
Tato geometrie silně motivuje exact no-hiding identity:
jestliže residue-side odd sektor skutečně odpovídá driftové komponentě a remainder je vázán
na odlišný sektor, pak je exaktní ortogonalita přirozená.
Ale samotný obraz ještě nedokazuje, že bank inner product přesně realizuje tuto sektorovou ortogonalitu.
To musí dodat analytická residue-side konstrukce.
\end{remark}

\subsection*{Lokalizovaná energie a topologický index}

\begin{definition}[Geometric core density]
V geometric residue-side modelu definujeme lokalizovanou energetickou hustotu
\[
\mathcal E(r,\theta):=|\psi(r,\theta)|^2.
\]
Řekneme, že mód je \emph{topologicky nenulový}, pokud jeho uzlová struktura a laloková
lokalizace vynucují nenulový winding/core index.
\end{definition}

\begin{proposition}[Non-collapse principle in the geometric model]
Je-li energetická hustota \(\mathcal E(r,\theta)\) lokalizovaná do stabilních laloků
a uzlová struktura nese nenulový winding, pak mód nemůže geometricky kolabovat do nulové amplitudy
bez ztráty své topologické identity.
\end{proposition}

\begin{remark}[Co to znamená pro M4]
To je přesně geometrický obsah coercivity locku:
nenulový topologický index má nést nenulovou koercivní konstantu.
V jazyce v23 to znamená, že identita
\[
c_4 \equiv G\,|\operatorname{Ind}_{\mathrm{core}}(\rho)|
\]
není libovolný odhad, ale odpovídá přímému geometrickému požadavku:
mód s nenulovým winding nesmí mít nulovou residue-side energii.
Stále však chybí analytický důkaz, že witness banka čte tento winding právě s touto konstantou.
\end{remark}

\subsection*{Geometric intertwining picture}

\begin{proposition}[Phase-locked transport in the geometric model]
Je-li sektorová struktura stabilní a fáze \(\Phi(r,\theta)\) se transportem zachovává
na úrovni dominantního laloku, pak model vykazuje phase-locked evoluci:
transport posouvá mód, ale nerozmazává jeho vnitřní paritní strukturu.
\end{proposition}

\begin{remark}
To je geometrická verze crown stability seal a resonant intertwining principu.
Model říká, že stabilita neznamená klid, ale zachování koherentního tvaru během transportu.
Jako strukturální podpora je to silné; jako důkaz M4/M5 to ale stále musí být přeloženo
do operatorové residue-side normy.
\end{remark}

\subsection*{Geometric seal versus proof layer}

\begin{remark}[What the images do prove, and what they do not]
Tyto geometrické modely dělají tři důležité věci:
\begin{enumerate}[label=\textbf{(\arabic*)}]
    \item dávají konkrétní residue-side tvar odd sektoru;
    \item dávají model nenulové lokalizované energie a uzlové struktury;
    \item dávají intuitivní mechaniku phase-locking a sektorové separace.
\end{enumerate}
Nedokazují však samy:
\[
Q_{\mathrm{cross}}(\rho,\mathrm{rest};\tau,L)=0,
\qquad
c_4 \equiv G\,|\operatorname{Ind}_{\mathrm{core}}(\rho)|
\]
v analytickém smyslu witness banky.
Tyto dvě identity zůstávají i ve v25 exact lock targets.
\end{remark}

\begin{remark}[Why this layer still matters]
Přesto je geometric residue-side seal důležitý, protože odstraňuje poslední dojem,
že residue-side lock targets jsou čistě formalistické.
Naopak:
\begin{enumerate}[label=\textbf{(\arabic*)}]
    \item odd sektor má konkrétní geometrický nosič,
    \item nenulový core index má konkrétní geometrický výraz,
    \item no-hiding a coercivity mají přímý vizuálně-modelový význam.
\end{enumerate}
Tím se v25 stává nejúplnější verzí programu: algebra, topologie, statistika, transport a geometrie
jsou vepsány do jediného lock-form manuscriptu.
\end{remark}



\section{Version 26: operator-valued residue-side addendum}

Tato sekce převádí geometric residue-side seal do explicitního operátorového jazyka.
Jejím cílem je ukázat, že čtyřlaločná struktura není jen obrázek, ale konečně-módový stav
v angular-momentum bázi, na který působí obecně nediagonální fázový operátor.
Tím se residue-side geometrie přepisuje do čisté spektrální formy bez ztráty struktury.

\subsection*{Angular-momentum Hilbertův prostor}

\begin{definition}[Angular sector Hilbert space]
Nechť
\[
\mathcal H_\theta:=L^2(S^1)
\]
je Hilbertův prostor čtvercově integrovatelných funkcí na kružnici.
Definujme operátor úhlového momentu
\[
L_z:=-i\partial_\theta.
\]
Pak Fourierovy módy
\[
e_m(\theta):=e^{im\theta},
\qquad m\in\mathbb Z,
\]
splňují
\[
L_z e_m = m e_m.
\]
\end{definition}

\begin{remark}
To je přesně operátorové čtení angular decomposition:
módy \(e^{im\theta}\) jsou vlastní stavy \(L_z\) a index \(m\) hraje roli spektrálního
angular momenta.
\end{remark}

\subsection*{Three-mode sector}

\begin{definition}[Three-mode seed state]
Použijeme identitu
\[
\cos(4\theta)=\frac12\bigl(e^{i4\theta}+e^{-i4\theta}\bigr).
\]
Proto modelovou amplitudovou část
\[
a+b\cos(4\theta)
\]
přepíšeme jako
\[
a+\frac b2 e^{i4\theta}+\frac b2 e^{-i4\theta}.
\]
Tedy seed angular state je
\[
\chi(\theta):=a+\frac b2 e^{i4\theta}+\frac b2 e^{-i4\theta}.
\]
\end{definition}

\begin{definition}[Projection onto the \(\{0,\pm4\}\)-sector]
Definujme ortogonální projekci
\[
\Pi_{0,\pm4}
=
|0\rangle\langle 0|+|4\rangle\langle 4|+|-4\rangle\langle -4|
\]
na podprostor
\[
\mathcal H_{0,\pm4}:=\mathrm{span}\{e^{i0\theta},e^{i4\theta},e^{-i4\theta}\}.
\]
\end{definition}

\begin{remark}
To je přesná verze toho, co bylo dříve jen naznačeno symbolem typu
\[
a+\frac b2(e^{i4\theta}+e^{-i4\theta}).
\]
Tento výraz sám o sobě není projekce; skutečná projekce je operátor \(\Pi_{0,\pm4}\).
\end{remark}

\subsection*{Residue-side wave state}

\begin{definition}[Operator-form residue-side state]
Nechť \(R(r)\) je radiální amplituda a nechť \(\Phi(r,\theta)\) je fáze.
Definujme unitární fázový operátor
\[
U_\Phi f(r,\theta):=e^{i\Phi(r,\theta)}f(r,\theta).
\]
Pak residue-side stav zapisujeme ve tvaru
\[
\psi(r,\theta)=U_\Phi\Bigl(R(r)\chi(\theta)\Bigr)
=
R(r)e^{i\Phi(r,\theta)}
\left(a+\frac b2 e^{i4\theta}+\frac b2 e^{-i4\theta}\right).
\]
\end{definition}

\begin{remark}
Jestliže \(\Phi=\Phi(r)\), pak \(U_\Phi\) je diagonální vzhledem k angular módům
a stav zůstává přesně v sektoru \(\{0,\pm4\}\).
Jestliže \(\Phi=\Phi(r,\theta)\), vzniká módové míchání.
Právě to je operátorová verze swirl / vortex struktury.
\end{remark}

\subsection*{Norma a Parsevalova identita}

\begin{proposition}[Three-mode norm identity]
Jestliže \(\Phi=\Phi(r)\), pak pro každý pevný \(r\) platí
\[
\int_0^{2\pi} |\psi(r,\theta)|^2\,d\theta
=
2\pi |R(r)|^2
\left(|a|^2+\frac{|b|^2}{2}\right).
\]
Tedy
\[
\|\psi\|^2
=
2\pi\int |R(r)|^2
\left(|a|^2+\frac{|b|^2}{2}\right)\,dr.
\]
\end{proposition}

\begin{proof}
Při \(\Phi=\Phi(r)\) je faktor \(e^{i\Phi(r)}\) globální fází v angular části a nemění normu.
Použijeme ortogonalitu Fourierových módů:
\[
\int_0^{2\pi} e^{im\theta}\overline{e^{in\theta}}\,d\theta
=
2\pi\delta_{mn}.
\]
Proto křížové členy mezi \(m=0\) a \(m=\pm4\), resp. mezi \(m=4\) a \(m=-4\), vymizí a zůstane
\[
2\pi|R(r)|^2\left(|a|^2+\left|\frac b2\right|^2+\left|\frac b2\right|^2\right)
=
2\pi|R(r)|^2\left(|a|^2+\frac{|b|^2}{2}\right).
\]
Integrací v \(r\) dostáváme druhý vzorec.
\end{proof}

\begin{remark}
Tato identita ukazuje, že residue-side energie v čistém \(\{0,\pm4\}\)-sektoru je nesena
koeficienty módů, nikoli jen vizuální kresbou laloků.
To je důležité pro přechod od geometric seal k operator-side norm.
\end{remark}

\subsection*{Módové míchání přes nediagonální fázi}

\begin{proposition}[Mode mixing under nontrivial phase]
Nechť
\[
e^{i\Phi(r,\theta)}=\sum_{k\in\mathbb Z}\alpha_k(r)e^{ik\theta}
\]
je Fourierův rozvoj unitárního fázového faktoru.
Pak výsledný residue-side stav má angular rozvoj
\[
\psi(r,\theta)
=
R(r)\sum_{n\in\mathbb Z}\widetilde c_n(r)e^{in\theta},
\]
kde koeficienty splňují konvoluční pravidlo
\[
\widetilde c_n(r)
=
a\,\alpha_n(r)+\frac b2\,\alpha_{n-4}(r)+\frac b2\,\alpha_{n+4}(r).
\]
\end{proposition}

\begin{proof}
Dosadíme Fourierův rozvoj \(e^{i\Phi(r,\theta)}\) do vztahu
\[
\psi(r,\theta)=
R(r)e^{i\Phi(r,\theta)}
\left(a+\frac b2 e^{i4\theta}+\frac b2 e^{-i4\theta}\right)
\]
a přeskupíme exponenty \(e^{i(n)\theta}\).
Koeficient u \(e^{in\theta}\) je právě
\[
a\,\alpha_n(r)+\frac b2\,\alpha_{n-4}(r)+\frac b2\,\alpha_{n+4}(r).
\]
\end{proof}

\begin{remark}
Tohle je přesný operátorový obsah tvrzení:
\[
\text{swirl} \Longleftrightarrow \text{fázový operátor je nediagonální v angular bázi.}
\]
Módové míchání není metafora, ale konvoluce spektrálních koeficientů.
\end{remark}

\subsection*{Residue-side Laurent interpretation}

\begin{definition}[Laurent-side symbol]
Položme
\[
z=e^{i\theta}.
\]
Pak angular seed symbol lze zapsat jako Laurentův polynom
\[
f(z)=a+\frac b2 z^4+\frac b2 z^{-4}.
\]
\end{definition}

\begin{remark}
Tento zápis není jen formální:
\begin{enumerate}[label=\textbf{(\arabic*)}]
    \item kladné a záporné mocniny kódují módové sektory \(m=\pm4\),
    \item člen \(z^{-4}\) je residue-side záporný mód,
    \item celý seed sektor je konečně-módový Laurent/Fourier objekt.
\end{enumerate}
Je ale důležité odlišit:
\[
\text{Laurent-side modal weight}
\quad\neq\quad
\text{classical complex residue}
\]
pokud není výslovně zavedena další konturová analytická struktura.
\end{remark}

\subsection*{Operator-side meaning for the bridge program}

\begin{remark}[Bridge relevance]
Operator-valued residue-side addendum posiluje program ve čtyřech bodech:
\begin{enumerate}[label=\textbf{(\arabic*)}]
    \item odd/parity sektor dostává konkrétní spektrální reprezentaci;
    \item nenulová laloková geometrie dostává konkrétní módovou normu;
    \item swirl / vortex struktura dostává přesný operátorový význam jako módové míchání;
    \item residue-side normalizace se dá číst zároveň geometricky, operatorově i Laurentovsky.
\end{enumerate}
Tím se v25 geometrická vrstva zvedá do operator-side jazyka bez ztráty struktury.
\end{remark}

\begin{remark}[Honest status]
Ani tato sekce sama ještě nedokazuje exact coercivity lock nebo exact no-hiding identity
v analytickém smyslu witness banky.
Co ale dává, je přesná operator-side forma residue-side modelu, ve které lze tyto dva locky
formulovat bez neurčitosti.
To je přesně správný krok od ``hezkého obrázku'' k ``reálné operátorové fyzice''.
\end{remark}


\section{Krátký verdict}

Celý program má tedy formu:
\[
\boxed{
\text{Closure} \Longrightarrow E_{\mathrm{ex}}(f)=0
}
\]
a současně
\[
\boxed{
\neg \mathrm{RH}
\Longrightarrow
\text{transportovaný drift}
\Longrightarrow
\text{Target Bridge Inequality}
\Longrightarrow
E_{\mathrm{ex}}(f)>0.
}
\]
Proto
\[
0\neq +\infty
\]
a existence mimokritického nulového bodu vede ke sporu.

\bigskip

\noindent
\textbf{Q.E.D. spine:}
\[
\text{closure}+\text{bridge}\Longrightarrow \mathrm{RH}.
\]

\begin{remark}
Tato verze obsahuje i Catalan parity route, corrected asymptotic firewall, crown stability seal, indistinguishability seal, generator distinction seal, exact residue-side lock targets, geometric residue-side seal a operator-valued residue-side addendum.
Tyto vrstvy výrazně zužují otevřené místo v bridge programu; zbytek je nyní přesně komprimován do dvou exact lock identities v residue-side normalizaci.
\end{remark}

\end{document}
